% Required Packages: \usepackage{graphicx}, \usepackage{float}

\begin{figure}[htbp]
  \centering
  \begin{subfigure}[b]{0.48\linewidth}
    \centering
    \zimg{rf_fig_0.jpg}{width=\linewidth,max height=0.4\textheight,keepaspectratio}{fg_0_14ct2}{rf_fig_0.jpg}
    \caption{(a) A representation of an unrolled graphene sheet along with the chiral vector definition (b) SWNT and MWNT}

  \end{subfigure}
  \hfill
  \begin{subfigure}[b]{0.48\linewidth}
    \centering
    \zimg{rf_fig_1.jpg}{width=\linewidth,max height=0.4\textheight,keepaspectratio}{fg_1_02on6}{rf_fig_1.jpg}
    \caption{(a) A representation of an unrolled graphene sheet along with the chiral vector definition (b) SWNT and MWNT}

  \end{subfigure}
  \caption{(a) A representation of an unrolled graphene sheet along with the chiral vector definition (b) SWNT}

  \label{fig:group_tfvnf}
\end{figure}

\begin{figure}[htbp]
  \centering
  \begin{subfigure}[b]{0.48\linewidth}
    \centering
    \zimg{rf_fig_2.png}{width=\linewidth,max height=0.4\textheight,keepaspectratio}{fg_0_4k3j2}{rf_fig_2.png}
    \caption{(c) CNTFET's three-dimensional structure (d)The CNTFET's top view with its parameters}

  \end{subfigure}
  \hfill
  \begin{subfigure}[b]{0.48\linewidth}
    \centering
    \zimg{rf_fig_3.png}{width=\linewidth,max height=0.4\textheight,keepaspectratio}{fg_1_71ay7}{rf_fig_3.png}
    \caption{(c) CNTFET's three-dimensional structure (d)The CNTFET's top view with its parameters}

  \end{subfigure}
  \caption{(c) CNTFET's three-dimensional structure (d)The CNTFET's top view with its parameters}

  \label{fig:group_yk4t5}
\end{figure}

\begin{figure}[htbp]
  \centering
  \begin{subfigure}[b]{0.23\linewidth}
    \centering
    \zimg{rf_fig_4.png}{width=\linewidth,max height=0.3\textheight,keepaspectratio}{fg_0_h3nhv}{rf_fig_4.png}
    \caption{(a) Pure CNTFET TC-OTA (b) NCNTFET-PMOS-TC-OTA (c) PCNTFET-NMOS-TC-OTA (d) Bulk CMOS-based TC-OTA}

  \end{subfigure}
  \hfill
  \begin{subfigure}[b]{0.23\linewidth}
    \centering
    \zimg{rf_fig_5.png}{width=\linewidth,max height=0.3\textheight,keepaspectratio}{fg_1_n9l25}{rf_fig_5.png}
    \caption{(a) Pure CNTFET TC-OTA (b) NCNTFET-PMOS-TC-OTA (c) PCNTFET-NMOS-TC-OTA (d) Bulk CMOS-based TC-OTA}

  \end{subfigure}
  \hfill
  \begin{subfigure}[b]{0.23\linewidth}
    \centering
    \zimg{rf_fig_6.png}{width=\linewidth,max height=0.3\textheight,keepaspectratio}{fg_2_rws0j}{rf_fig_6.png}
    \caption{(a) Pure CNTFET TC-OTA (b) NCNTFET-PMOS-TC-OTA (c) PCNTFET-NMOS-TC-OTA (d) Bulk CMOS-based TC-OTA}

  \end{subfigure}
  \hfill
  \begin{subfigure}[b]{0.23\linewidth}
    \centering
    \zimg{rf_fig_7.png}{width=\linewidth,max height=0.3\textheight,keepaspectratio}{fg_3_amro4}{rf_fig_7.png}
    \caption{(a) Pure CNTFET TC-OTA (b) NCNTFET-PMOS-TC-OTA (c) PCNTFET-NMOS-TC-OTA (d) Bulk CMOS-based TC-OTA}

  \end{subfigure}
  \caption{(a) Pure CNTFET TC-OTA (b) NCNTFET-PMOS-TC-OTA (c) PCNTFET-NMOS-TC-OTA (d) Bulk CMOS-based TC-OTA}

  \label{fig:group_pgg2f}
\end{figure}

\begin{figure}[htbp]
  \centering
  \begin{subfigure}[b]{0.48\linewidth}
    \centering
    \zimg{rf_fig_8.png}{width=\linewidth,max height=0.4\textheight,keepaspectratio}{fg_0_r1w9m}{rf_fig_8.png}
    \caption{(b) shows that the output resistance in all TC-OTAs diminishes as N grows. Since an increase in N is equivalent to an increase...}

  \end{subfigure}
  \hfill
  \begin{subfigure}[b]{0.48\linewidth}
    \centering
    \zimg{rf_fig_9.png}{width=\linewidth,max height=0.4\textheight,keepaspectratio}{fg_1_imi8a}{rf_fig_9.png}
    \caption{(b) shows that the output resistance in all TC-OTAs diminishes as N grows. Since an increase in N is equivalent to an increase...}

  \end{subfigure}
  \caption{(b) shows that the output resistance in all TC-OTAs diminishes as N grows. Since an increase in N is equivalent to an increase in a transistor's width W, output resistance (ROUT) must drop as N increases. However, TC-OTAs have a greater ROUT due to the cascoding effect. Figure 3(c) shows how N affects the average power consumption of CNT-based TC-OTAs. It is evident that average power grows along with N. This is explained by the fact that CNTFETs' driving capacity has increased due to the increase in N. However, because to the exceptionally low CNT channel resistance, 1D ballistic transport, reduced parasitic capacitance, and reduced leakage, a pure CNT-based TC-OTA has a much lower power dissipation \cite{ref18,ref35}. When comparing the two hybrid TC-OTAs, the NCNTFET-PMOS-TC-OTA uses more power than the PCNTFET-NMOS-TC-OTA because the former has more nCNTFETs (14 nCNTFETs) than the latter (12 pCNTFETs). The bandwidth is greatly enhanced in CNT-based TC-OTAs, as shown in Figure 3(d), with the amount of CNTs having the biggest impact on pure CNT-TC-OTA. It is more important for PCNTFET-NMOS-TC-OTA circuits than NCNTFET-PMOS-TC-OTA circuits, though, and grows linearly \cite{ref6,ref19}.}

  \label{fig:group_67bqx}
\end{figure}

\begin{figure}[htbp]
  \centering
  \begin{subfigure}[b]{0.48\linewidth}
    \centering
    \zimg{rf_fig_10.png}{width=\linewidth,max height=0.4\textheight,keepaspectratio}{fg_0_puhb4}{rf_fig_10.png}
    \caption{Impact of changing the number of CNTs in the proposed CNT-based TC-OTAs on (a) DC gain (b) output resistance, (c) average power, and (d) bandwidth}

  \end{subfigure}
  \hfill
  \begin{subfigure}[b]{0.48\linewidth}
    \centering
    \zimg{rf_fig_11.png}{width=\linewidth,max height=0.4\textheight,keepaspectratio}{fg_1_0dbak}{rf_fig_11.png}
    \caption{Impact of changing the number of CNTs in the proposed CNT-based TC-OTAs on (a) DC gain (b) output resistance, (c) average power, and (d) bandwidth}

  \end{subfigure}
  \caption{Impact of changing the number of CNTs in the proposed CNT-based TC-OTAs on (a) DC gain (b) output resistance, (c) average power, and (d) bandwidth}

  \label{fig:group_hyrsd}
\end{figure}

\begin{figure}[htbp]
  \centering
  \begin{subfigure}[b]{0.48\linewidth}
    \centering
    \zimg{rf_fig_12.png}{width=\linewidth,max height=0.4\textheight,keepaspectratio}{fg_0_8hndm}{rf_fig_12.png}
    \caption{Impact of changing the CNT diameter in the proposed CNT-based TC-OTAs on (a) DC gain (b) output resistance, (c) average power, and (d) bandwidth}

  \end{subfigure}
  \hfill
  \begin{subfigure}[b]{0.48\linewidth}
    \centering
    \zimg{rf_fig_13.png}{width=\linewidth,max height=0.4\textheight,keepaspectratio}{fg_1_nsem0}{rf_fig_13.png}
    \caption{Impact of changing the CNT diameter in the proposed CNT-based TC-OTAs on (a) DC gain (b) output resistance, (c) average power, and (d) bandwidth}

  \end{subfigure}
  \caption{Impact of changing the CNT diameter in the proposed CNT-based TC-OTAs on (a) DC gain (b) output resistance, (c) average power, and (d) bandwidth}

  \label{fig:group_wxw2q}
\end{figure}

\begin{figure}[htbp]
  \centering
  \begin{subfigure}[b]{0.48\linewidth}
    \centering
    \zimg{rf_fig_14.png}{width=\linewidth,max height=0.4\textheight,keepaspectratio}{fg_0_uk4l7}{rf_fig_14.png}
    \caption{Impact of changing the CNT diameter in the proposed CNT-based TC-OTAs on (a) DC gain (b) output resistance, (c) average power, and (d) bandwidth}

  \end{subfigure}
  \hfill
  \begin{subfigure}[b]{0.48\linewidth}
    \centering
    \zimg{rf_fig_15.png}{width=\linewidth,max height=0.4\textheight,keepaspectratio}{fg_1_2wdaz}{rf_fig_15.png}
    \caption{Impact of changing the CNT diameter in the proposed CNT-based TC-OTAs on (a) DC gain (b) output resistance, (c) average power, and (d) bandwidth}

  \end{subfigure}
  \caption{Impact of changing the CNT diameter in the proposed CNT-based TC-OTAs on (a) DC gain (b) output resistance, (c) average power, and (d) bandwidth}

  \label{fig:group_7084h}
\end{figure}

\begin{figure}[htbp]
  \centering
  \begin{subfigure}[b]{0.48\linewidth}
    \centering
    \zimg{rf_fig_16.png}{width=\linewidth,max height=0.4\textheight,keepaspectratio}{fg_0_9vg48}{rf_fig_16.png}
  \end{subfigure}
  \hfill
  \begin{subfigure}[b]{0.48\linewidth}
    \centering
    \zimg{rf_fig_17.png}{width=\linewidth,max height=0.4\textheight,keepaspectratio}{fg_1_j6sh5}{rf_fig_17.png}
    \caption{Impact of changing the CNT pitch in the proposed CNT-based TC-OTAs on (a) DC gain (b) output resistance, (c) average power, and (d) bandwidth}

  \end{subfigure}
  \caption{Impact of changing the CNT pitch in the proposed CNT-based TC-OTAs on (a) DC gain (b) output resistance, (c) average power, and (d) bandwidth}

  \label{fig:group_2ewis}
\end{figure}

\begin{figure}[htbp]
  \centering
  \begin{subfigure}[b]{0.48\linewidth}
    \centering
    \zimg{rf_fig_18.png}{width=\linewidth,max height=0.4\textheight,keepaspectratio}{fg_0_rct90}{rf_fig_18.png}
    \caption{Impact of changing the CNT pitch in the proposed CNT-based TC-OTAs on (a) DC gain (b) output resistance, (c) average power, and (d) bandwidth}

  \end{subfigure}
  \hfill
  \begin{subfigure}[b]{0.48\linewidth}
    \centering
    \zimg{rf_fig_19.png}{width=\linewidth,max height=0.4\textheight,keepaspectratio}{fg_1_sknmm}{rf_fig_19.png}
    \caption{Impact of changing the CNT pitch in the proposed CNT-based TC-OTAs on (a) DC gain (b) output resistance, (c) average power, and (d) bandwidth}

  \end{subfigure}
  \caption{Impact of changing the CNT pitch in the proposed CNT-based TC-OTAs on (a) DC gain (b) output resistance, (c) average power, and (d) bandwidth}

  \label{fig:group_eb2xz}
\end{figure}

\begin{figure}[htbp]
  \centering
  \begin{subfigure}[b]{0.48\linewidth}
    \centering
    \zimg{rf_fig_20.png}{width=\linewidth,max height=0.4\textheight,keepaspectratio}{fg_0_12qgj}{rf_fig_20.png}
    \caption{(a) shows the behavior of various operational transconductance amplifier (OTA) topologies' DC Gain as the transistors' oxide...}

  \end{subfigure}
  \hfill
  \begin{subfigure}[b]{0.48\linewidth}
    \centering
    \zimg{rf_fig_21.png}{width=\linewidth,max height=0.4\textheight,keepaspectratio}{fg_1_1dbze}{rf_fig_21.png}
    \caption{(a) shows the behavior of various operational transconductance amplifier (OTA) topologies' DC Gain as the transistors' oxide...}

  \end{subfigure}
  \caption{(a) shows the behavior of various operational transconductance amplifier (OTA) topologies' DC Gain as the transistors' oxide thickness (TOX) varies between 2 nm and 5 nm. Gate capacitance and drain current are often significantly impacted by the thickness of the gate oxide. Although it might result in leakage, a thinner oxide often improves gate control. The transconductance (gmn) and output resistance (ROCASN) of the transistors, which determine DC gain (AV=gmn* ROUT), seem to be leveling out or being controlled by other dominating elements in the OTA within this particular range, according to the flat lines \cite{ref7,ref39}.}

  \label{fig:group_y7snq}
\end{figure}

\begin{figure}[htbp]
  \centering
  \begin{subfigure}[b]{0.48\linewidth}
    \centering
    \zimg{rf_fig_22.png}{width=\linewidth,max height=0.4\textheight,keepaspectratio}{fg_0_zzqrt}{rf_fig_22.png}
    \caption{Impact of changing the oxide thickness in the proposed CNT-based TC-OTAs on (a) DC gain (b) output resistance, (c) average power, and (d) bandwidth}

  \end{subfigure}
  \hfill
  \begin{subfigure}[b]{0.48\linewidth}
    \centering
    \zimg{rf_fig_23.png}{width=\linewidth,max height=0.4\textheight,keepaspectratio}{fg_1_4wmie}{rf_fig_23.png}
    \caption{Impact of changing the oxide thickness in the proposed CNT-based TC-OTAs on (a) DC gain (b) output resistance, (c) average power, and (d) bandwidth}

  \end{subfigure}
  \caption{Impact of changing the oxide thickness in the proposed CNT-based TC-OTAs on (a) DC gain (b) output resistance, (c) average power, and (d) bandwidth}

  \label{fig:group_jnhxf}
\end{figure}

\begin{figure}[htbp]
  \centering
  \begin{subfigure}[b]{0.48\linewidth}
    \centering
    \zimg{rf_fig_24.png}{width=\linewidth,max height=0.4\textheight,keepaspectratio}{fg_0_ac56c}{rf_fig_24.png}
    \caption{Impact of changing the dielectric constant in the proposed CNT-based TC-OTAs on (a) DC gain (b) output resistance, (c) average...}

  \end{subfigure}
  \hfill
  \begin{subfigure}[b]{0.48\linewidth}
    \centering
    \zimg{rf_fig_25.png}{width=\linewidth,max height=0.4\textheight,keepaspectratio}{fg_1_8y7uh}{rf_fig_25.png}
    \caption{Impact of changing the dielectric constant in the proposed CNT-based TC-OTAs on (a) DC gain (b) output resistance, (c) average...}

  \end{subfigure}
  \caption{Impact of changing the dielectric constant in the proposed CNT-based TC-OTAs on (a) DC gain (b) output resistance, (c) average power, and (d) bandwidth}

  \label{fig:group_m4kt3}
\end{figure}

\begin{figure}[htbp]
  \centering
  \begin{subfigure}[b]{0.48\linewidth}
    \centering
    \zimg{rf_fig_26.png}{width=\linewidth,max height=0.4\textheight,keepaspectratio}{fg_0_g4rwk}{rf_fig_26.png}
    \caption{Impact of changing the dielectric constant in the proposed CNT-based TC-OTAs on (a) DC gain (b) output resistance, (c) average...}

  \end{subfigure}
  \hfill
  \begin{subfigure}[b]{0.48\linewidth}
    \centering
    \zimg{rf_fig_27.png}{width=\linewidth,max height=0.4\textheight,keepaspectratio}{fg_1_zgxtj}{rf_fig_27.png}
    \caption{Impact of changing the dielectric constant in the proposed CNT-based TC-OTAs on (a) DC gain (b) output resistance, (c) average...}

  \end{subfigure}
  \caption{Impact of changing the dielectric constant in the proposed CNT-based TC-OTAs on (a) DC gain (b) output resistance, (c) average power, and (d) bandwidth}

  \label{fig:group_r9a9f}
\end{figure}

\begin{figure}[htbp]
  \centering
  \begin{subfigure}[b]{0.48\linewidth}
    \centering
    \zimg{rf_fig_28.png}{width=\linewidth,max height=0.4\textheight,keepaspectratio}{fg_0_9sw0q}{rf_fig_28.png}
    \caption{Frequency response of different TC-OTAs.}

  \end{subfigure}
  \hfill
  \begin{subfigure}[b]{0.48\linewidth}
    \centering
    \zimg{rf_fig_29.png}{width=\linewidth,max height=0.4\textheight,keepaspectratio}{fg_1_vdro3}{rf_fig_29.png}
    \caption{Frequency response of different TC-OTAs.}

  \end{subfigure}
  \caption{Frequency response of different TC-OTAs.}

  \label{fig:group_ieujb}
\end{figure}

\begin{figure}[htbp]
  \centering
  \begin{subfigure}[b]{0.48\linewidth}
    \centering
    \zimg{rf_fig_30.png}{width=\linewidth,max height=0.4\textheight,keepaspectratio}{fg_0_khdm7}{rf_fig_30.png}
    \caption{Frequency response of different TC-OTAs.}

  \end{subfigure}
  \hfill
  \begin{subfigure}[b]{0.48\linewidth}
    \centering
    \zimg{rf_fig_31.png}{width=\linewidth,max height=0.4\textheight,keepaspectratio}{fg_1_sei6e}{rf_fig_31.png}
    \caption{Frequency response of different TC-OTAs.}

  \end{subfigure}
  \caption{Frequency response of different TC-OTAs.}

  \label{fig:group_ncpui}
\end{figure}

\begin{figure}[htbp]
  \centering
  \begin{subfigure}[b]{0.48\linewidth}
    \centering
    \zimg{rf_fig_32.png}{width=\linewidth,max height=0.4\textheight,keepaspectratio}{fg_0_g99vu}{rf_fig_32.png}
    \caption{(a) Proposed CNT-OTA-Based BPF (b) All four OTAs' frequency responses for the specified BPF}

  \end{subfigure}
  \hfill
  \begin{subfigure}[b]{0.48\linewidth}
    \centering
    \zimg{rf_fig_33.png}{width=\linewidth,max height=0.4\textheight,keepaspectratio}{fg_1_vj307}{rf_fig_33.png}
    \caption{(a) Proposed CNT-OTA-Based BPF (b) All four OTAs' frequency responses for the specified BPF}

  \end{subfigure}
  \caption{(a) Proposed CNT-OTA-Based BPF (b) All four OTAs' frequency responses for the specified BPF}

  \label{fig:group_7xrdu}
\end{figure}
